\scanpage{64}
% manuscript line 1
En effet, il suffit, dans (4), de \unsure{s'écrire|s'aider de} l'inégalité de
Cauchy-Schwartz pour le premier mb, et d'utiliser (5).
Ce corollaire prouve que le coefficient $k_n$ qui intervient
dans Chap. 1, § 2, N° 2, prop. 1 est égal à 1 pour \deleted{l'} \inserted{un} espace de
Hilbert, (\deleted{\unsure{compte tenu}} \inserted{compte tenu} de l'identification \deleted{\unsure{du}} dual $E'$ de $E$ à $\bar E$)

% manuscript line 6
Soient maintenant $E,F$ deux espaces de Hilbert, et
soient $u_1,\ldots u_n$ $n$ applications linéaires continues
de $E$ dans $F$, considérons leur produit extérieur
$u_1\wedge\cdots\wedge u_n$ (Chap 1, § 1, N° 1). On a

% manuscript line 10
\[
u_1\wedge\cdots\wedge u_n=\deleted{\unsure{n!|1/n!}}\,\varphi_n\circ(u_1\otimes\cdots\otimes u_n)\circ a_n \tag*{(8)}
\]

% manuscript line 11
où $a_n$ est l'isomorphisme \deleted{\unsure{naturel}} naturel de
$\overset{n}{\bigwedge}E$ dans $\overset{n}{\otimes}E$, et $\varphi_n$ l'homomorphisme naturel de
$\overset{n}{\otimes}F$ \deleted{\unsure{dans}} \unsure{sur} $\overset{n}{\bigwedge}F$. Comme $\frac1{\sqrt{n!}}a_n$ est \unsure{$=$} un
isomorphisme de structures préhilbertiennes, et $\sqrt{n!}\,\varphi_n$ en est
de norme $\leq1$ pour les normes hilbertiennes mises
sur $\overset{n}{\otimes}F$ et $\overset{n}{\bigwedge}F$, on voit que

% manuscript line 17
\[
\|u_1\wedge\cdots\wedge u_n\|=\|\sqrt{n!}\,\varphi_n\circ(u_1\otimes\cdots\otimes u_n)\circ\frac1{\sqrt{n!}}a_n\|\leq\|u_1\otimes\cdots\otimes u_n\|,
\]

% manuscript line 18
d'où, en vertu de (2) :

% manuscript line 19
\[
\|u_1\wedge\cdots\wedge u_n\|\leq\|u_1\|\cdots\|u_n\|, \tag*{(9)}
\]

% manuscript line 20
Par suite, $u_1\wedge\cdots\wedge u_n$ peut être étendu en
une application linéaire continue de $\overset{n}{\bigwedge}\vphantom{\bigwedge}^{(2)}E$ dans
$\overset{n}{\bigwedge}\vphantom{\bigwedge}^{(2)}F$, \inserted{notée encore $u_1\wedge\cdots\wedge u_n$,} dont la norme satisfait à (9). Bien entendu,
les règles de calcul usuelles subsis\deleted{\unsure{tent}}tent pour les
\deleted{\unsure{produits}} $u_1\wedge\cdots\wedge u_n$. \inserted{(prolongés par continuité)} En particulier, si $u_i$ est un

% manuscript line 25
% archive sheet number: 65
